\section{Introduction}
\label{sec:intro}

The electrocardiogram (ECG) is the primary non-invasive tool for assessing cardiac electrophysiology, and cardiovascular disease remains the leading cause of death worldwide \cite{clifford2006}. Continuous monitoring is moving from the clinic to smartwatches, Holter patches, and chest straps, where the recordings are contaminated by baseline wander (BW) from respiration, muscle artifacts (MA) from skeletal activity, and electrode motion (EM) at the skin--sensor interface \cite{clifford2006,jin2024}. Because MA and EM overlap the spectral band of the P wave, QRS complex, and T wave, this contamination cannot be removed by frequency selection alone without damaging the waveform on which clinicians and algorithms rely.

Classical remedies include bandpass and adaptive filtering, Kalman smoothing, empirical mode decomposition, and wavelet thresholding \cite{clifford2006,goldberger2000}. Learning-based methods replace hand-tuned rules with end-to-end mappings from noisy to clean segments, such as denoising autoencoders, fully convolutional networks, and U-Net variants \cite{vincent2008,chiang2019,xiong2016,ronneberger2015}, together with dilated multi-branch filters such as DeepFilter \cite{romero2021}. More recently, hybrid networks insert the discrete wavelet transform (DWT) into the network to obtain multiresolution features without pooling loss \cite{jin2024,nguyen2026dwse,nguyen2026liwave}.

\begin{figure*}[!t]
\centering
\begin{adjustbox}{max width=0.98\textwidth}
\begin{tikzpicture}[
  node distance=4mm,
  blk/.style={draw, rectangle, rounded corners=2pt, minimum height=7mm, minimum width=1.4cm, align=center, font=\footnotesize, fill=white, line width=0.6pt},
  den/.style={blk, fill=green!8, draw=green!60!black},
  cls/.style={blk, fill=blue!8, draw=blue!60!black},
  qnt/.style={blk, fill=violet!10, draw=violet!70!black},
  lossn/.style={blk, fill=orange!10, draw=orange!70!black},
  arr/.style={-{Stealth[length=2mm]}, thick},
  darr/.style={-{Stealth[length=2mm]}, thick, dashed},
  grp/.style={draw=black!60, dashed, rounded corners=4pt, inner sep=3.5mm, font=\scriptsize\bfseries}
]

% ==================== STAGE 1 ====================
\node[blk] (n1) {Noisy Segment $x[n]$\\($8192$ samples, $\mathrm{mV}$)};
\node[den, right=6mm of n1] (n2) {\textbf{\HSL}\\\textit{Wavelet Denoiser}};
\node[blk, right=6mm of n2] (n3) {Denoised Output\\$\hat{y}[n] = x + \Delta$};
\node[lossn, right=8mm of n3] (ls) {Mixed Loss $\mathcal{L}_{\mathrm{mixed}}$\\$0.85\mathcal{L}_{\mathrm{MSE}} + 0.15\mathcal{L}_{L_1}$};
\node[blk, below=4mm of ls] (n4) {Clean Target $x_0[n]$\\($8192$ samples, $\mathrm{mV}$)};

\draw[arr] (n1) -- (n2);
\draw[arr] (n2) -- (n3);
\draw[arr] (n3) -- (ls);
\draw[arr] (n4) -- (ls);

\begin{scope}[on background layer]
\node[grp, fit=(n1)(n2)(n3)(ls)(n4), label={[font=\footnotesize\bfseries, anchor=south west]north west:\textbf{Stage 1: Upstream Denoising (\HSL{} Trained Here)}}] (g1) {};
\end{scope}

% ==================== STAGE 2 ====================
\node[blk, below=26mm of n1] (m1) {Segmented Beat $x_d$\\($256$ samples, R-peak)};
\node[den, right=5mm of m1] (m2) {\textbf{\HSL}\\(Weights Frozen)};
\node[cls, right=5mm of m2] (m3) {Adaptive\\Z-Score Norm};
\node[cls, right=5mm of m3] (m4) {1D ResNet\\Morphology (128D)};
\node[cls, right=5mm of m4] (m5) {Feature Fusion\\(160D)};
\node[qnt, right=5mm of m5] (m6) {Classification Head\\(Classical MLP / VQC)};
\node[blk, right=5mm of m6] (m7) {Diagnostic Output\\\textbf{N / S / V}};

% RR Context Branch
\node[cls, below=6mm of m3] (r1) {6D RR-Intervals\\($f_1, \dots, f_6$)};
\node[cls, right=5mm of r1] (r2) {BN-Linear-GELU\\Rhythm Pathway (32D)};

\draw[arr] (m1) -- (m2);
\draw[arr] (m2) -- (m3);
\draw[arr] (m3) -- (m4);
\draw[arr] (m4) -- (m5);
\draw[arr] (m5) -- (m6);
\draw[arr] (m6) -- (m7);

\draw[arr] (r1) -- (r2);
\draw[arr] (r2.east) -| (m5.south);

% Frozen weights arrow
\draw[darr] (n2.south) -- ++(0,-7mm) -| node[pos=0.25, above, font=\scriptsize\bfseries, text=black!80]{Weights Frozen \& Reused} (m2.north);

\begin{scope}[on background layer]
\node[grp, fit=(m1)(m2)(m7)(r1)(r2), label={[font=\footnotesize\bfseries, anchor=south west]north west:\textbf{Stage 2: Downstream Heartbeat Arrhythmia Classification (\RRN{} Trained on DS1/DS2)}}] (g2) {};
\end{scope}

\end{tikzpicture}
\end{adjustbox}
\caption{End-to-end framework for wearable ECG analysis. \textbf{Stage 1 (Top)}: \HSL{} is trained on 8192-sample raw-scale continuous segments under single and multi-source NSTDB noise with a physics-preserving mixed objective. \textbf{Stage 2 (Bottom)}: The frozen denoiser purifies segmented 256-sample cardiac beats, processed by \RRN{} alongside an independent 6-dimensional RR interval timing pathway. Dual interchangeable heads (Classical MLP and 6-qubit Hybrid VQC) perform final 3-class AAMI arrhythmia classification (N, S, V) on the 49{,}298-beat inter-patient DS2 test set.}
\label{fig:pipeline}
\end{figure*}

These approaches still share three limitations. Handcrafted thresholds and static filters cannot adapt to non-stationary noise and often erase low-amplitude P and T waves. Networks trained with a pure squared-error objective reduce the overall error by flattening the steep R-peak, and most wavelet hybrids process the high-frequency detail subbands far more coarsely than the approximation branch. Finally, denoisers are usually judged by SNR-type metrics alone, so it remains unclear whether a higher SNR preserves the morphology required for arrhythmia detection on unseen patients \cite{dechazal2004}.

Our research direction is therefore a lightweight wavelet denoiser that explicitly refines the high-frequency detail subbands, is optimized with an objective that protects the R-peak amplitude, and is evaluated along two separate axes: the quality of the restored signal and the detection capability of a downstream classifier under the inter-patient protocol (Fig.~\ref{fig:pipeline}). The main contributions of this work are as follows.
\begin{itemize}
    \item \textbf{\HSL{} denoiser:} a lightweight symmetric two-level Haar wavelet network built from a single residual block type with squeeze-and-excitation, in which the detail subbands are refined by dedicated skip blocks before inverse wavelet reconstruction.
    \item \textbf{Fiducial-preserving training:} a mixed squared- and absolute-error objective that limits R-peak blunting, assessed through clinical fiducial measurements in addition to SNR-type metrics.
    \item \textbf{Two-part evaluation with downstream detection:} a signal-quality assessment against five denoising baselines and an inter-patient arrhythmia detection assessment with \RRN{}, using classical and 6-qubit hybrid quantum heads.
\end{itemize}
